\section*{第 \ref{ch:g10:lewis-and-shape} 章——路易斯结构与分子的形状}

\begin{solution}{exo:g10:lewis-and-shape:1}
由两个\omterm{def:g7:atoms-and-molecules:atom}{原子}共用的一对\omterm{def:g9:inside-the-atom:nucleus}{电子}。一对\omterm{def:g10:electron-shells:valence}{价电子}
仅属于一个\omterm{def:g7:atoms-and-molecules:atom}{原子}、未被共用。
\end{solution}

\begin{solution}{exo:g10:lewis-and-shape:2}
H：1（比\omterm{def:g10:electron-shells:octet}{二隅规则}少 1 个\omterm{def:g9:inside-the-atom:nucleus}{电子}）。C：4，N：3，O：2（分别比\omterm{def:g10:electron-shells:octet}{八隅规则}少
4、3 和 2 个\omterm{def:g9:inside-the-atom:nucleus}{电子}）。
\end{solution}

\begin{solution}{exo:g10:lewis-and-shape:3}
\chemfig{H-\charge{0=\:,90=\:,270=\:}{Cl}}：氯\omterm{def:g7:atoms-and-molecules:atom}{原子}有一对共用\omterm{def:g9:inside-the-atom:nucleus}{电子}对和三
对\omterm{def:g10:lewis-and-shape:lone-pair}{孤电子对}，共 8 个\omterm{def:g9:inside-the-atom:nucleus}{电子}。
\end{solution}

\begin{solution}{exo:g10:lewis-and-shape:4}
$5 + 3 \times 1 = 8$ 个\omterm{def:g9:inside-the-atom:nucleus}{电子}，4 对（三个键，一对\omterm{def:g10:lewis-and-shape:lone-pair}{孤电子对}）。
\end{solution}

\begin{solution}{exo:g10:lewis-and-shape:5}
\ce{CH4}：四面体形。\ce{NH3}：三角锥形。\ce{H2O}：角形。\ce{CO2}：
直线形。
\end{solution}

\begin{solution}{exo:g10:lewis-and-shape:6}
\chemfig{H-\charge{90=\:,270=\:}{S}-H}：与水相似，硫周围有四个基团，
其中两对\omterm{def:g10:lewis-and-shape:lone-pair}{孤电子对}，为角形。
\end{solution}

\begin{solution}{exo:g10:lewis-and-shape:7}
$2 \times 5 = 10$ 个\omterm{def:g9:inside-the-atom:nucleus}{电子}，5 对。一个单键和 4 对\omterm{def:g10:lewis-and-shape:lone-pair}{孤电子对}，使每个氮\omterm{def:g7:atoms-and-molecules:atom}{原子}
分别只剩 6 个\omterm{def:g9:inside-the-atom:nucleus}{电子}；将两对\omterm{def:g10:lewis-and-shape:lone-pair}{孤电子对}变为共用\omterm{def:g9:inside-the-atom:nucleus}{电子}对，
得到\omterm{def:g10:lewis-and-shape:multiple-bond}{三键}，每个\omterm{def:g7:atoms-and-molecules:atom}{原子}上各有一对\omterm{def:g10:lewis-and-shape:lone-pair}{孤电子对}：
\chemfig{\charge{180=\:}{N}~\charge{0=\:}{N}}。
\end{solution}

\begin{solution}{exo:g10:lewis-and-shape:8}
\chemfig{H-C(-[2]H)(-[6]H)-\charge{90=\:,270=\:}{O}-H}：碳周围有四个键，
为四面体形；氧周围有两个键和两对\omterm{def:g10:lewis-and-shape:lone-pair}{孤电子对}，为角形。
\end{solution}

\begin{solution}{exo:g10:lewis-and-shape:9}
\ce{CO2} 中，碳有两个基团（两个\omterm{def:g10:lewis-and-shape:multiple-bond}{双键}），没有\omterm{def:g10:lewis-and-shape:lone-pair}{孤电子对}：它们
指向相反方向。\ce{H2O} 中，氧有四个基团（两个键、两对\omterm{def:g10:lewis-and-shape:lone-pair}{孤电子对}），
呈四面体排布：两个键之间形成
一定夹角。
\end{solution}

\begin{solution}{exo:g10:lewis-and-shape:10}
\chemfig{C(-[3]H)(-[5]H)=C(-[1]H)-[7]H}：每个碳周围有三个基团，
即一个\omterm{def:g10:lewis-and-shape:multiple-bond}{双键}和两个单键，呈平面三角形，键角接近
\ang{120}。
\end{solution}

\begin{solution}{exo:g10:lewis-and-shape:11}
碳位于中心，与四个氯\omterm{def:g7:atoms-and-molecules:atom}{原子}各形成一个单键，
每个氯\omterm{def:g7:atoms-and-molecules:atom}{原子}上有三对\omterm{def:g10:lewis-and-shape:lone-pair}{孤电子对}：碳周围有四个基团，为四面体形。
\end{solution}

\begin{solution}{exo:g10:lewis-and-shape:12}
乙醇 \chemfig{H-C(-[2]H)(-[6]H)-C(-[2]H)(-[6]H)-\charge{90=\:,270=\:}{O}-H}，以及
甲醚，
\chemfig{H-C(-[2]H)(-[6]H)-\charge{90=\:,270=\:}{O}-C(-[2]H)(-[6]H)-H}。
\end{solution}

\begin{solution}{exo:g10:lewis-and-shape:13}
\omterm{def:g10:lewis-and-shape:lone-pair}{孤电子对}占据的空间比\omterm{def:g10:lewis-and-shape:lone-pair}{成键电子对}略大，会使
键之间的距离挤近。氨有一对\omterm{def:g10:lewis-and-shape:lone-pair}{孤电子对}，水有两对：键角
由甲烷到氨再到水逐渐减小。
\end{solution}

\begin{solution}{exo:g10:lewis-and-shape:14}
\chemfig{H-C~\charge{0=\:}{N}}：碳周围有两个基团，为直线形。
\chemfig{H-C(-[6]H)=\charge{45=\:,315=\:}{O}}：碳周围有三个基团，呈平面
三角形。
\end{solution}

\begin{solution}{exo:g10:lewis-and-shape:15}
$3 \times 6 = 18$ 个\omterm{def:g9:inside-the-atom:nucleus}{电子}，9 对：
\chemfig{\charge{180=\:,270=\:}{O}=\charge{90=\:}{O}-\charge{0=\:,90=\:,270=\:}{O}}。中心氧\omterm{def:g7:atoms-and-molecules:atom}{原子}
有三个基团（一个\omterm{def:g10:lewis-and-shape:multiple-bond}{双键}、一个单键、一对\omterm{def:g10:lewis-and-shape:lone-pair}{孤电子对}），
\omterm{def:g7:atoms-and-molecules:molecule}{分子}为角形。
\end{solution}

\begin{solution}{pb:g10:lewis-and-shape:1}
\textbf{1.} \chemfig{H-C(-[2]H)(-[6]H)-H}，\chemfig{H-\charge{90=\:}{N}(-[6]H)-H}，
\chemfig{H-\charge{90=\:,270=\:}{O}-H}。

\textbf{2.} 每个\omterm{def:g7:atoms-and-molecules:molecule}{分子}中都有四个基团。

\textbf{3.} 甲烷无\omterm{def:g10:lewis-and-shape:lone-pair}{孤电子对}，氨有一对，水有两对。

\textbf{4.} 碳形成四个键；氧形成两个键（其另外两个基团
是\omterm{def:g10:lewis-and-shape:lone-pair}{孤电子对}，小棍模型无法表示）。

\textbf{5.} \chemfig{H-C(-[2]H)(-[6]H)-C(-[2]H)(-[6]H)-H}：一根小棍连接
两个碳球，共需六个氢球。

\textbf{6.} 四面体形，三角锥形，角形。

\textbf{7.} \omterm{def:g10:lewis-and-shape:lone-pair}{孤电子对}比\omterm{def:g10:lewis-and-shape:lone-pair}{成键电子对}占据更多空间，会挤近
键之间的距离；水有两对\omterm{def:g10:lewis-and-shape:lone-pair}{孤电子对}，键角最小。

\textbf{8.} 碳球上的孔之间的夹角为 \ang{109.5}；水的真实键角为
\ang{104.5}：氧球需要在 \ang{104.5} 处钻出自己的两个孔。

\textbf{9.} \ang{180}：碳周围两个基团指向
相反方向，\omterm{def:g7:atoms-and-molecules:molecule}{分子}为直线形。

\textbf{10.} 不在同一条棱上的两个顶点，是同一个面的对角顶点：
它们的距离为面对角线，$\sqrt{a^2 + a^2} = a\sqrt{2}$。

\textbf{11.} 中心位于体对角线的中点：
$a\sqrt{3}/2$。

\textbf{12.} H--H：$\sqrt{2} \approx \qty{1.41}{cm}$；C--H：
$\sqrt{3}/2 \approx \qty{0.87}{cm}$。

\textbf{13.} 每个氢位于立方体的一个顶点，碳位于其
中心：每个顶点到中心的距离都相同。

\textbf{14.} 三角形 C\,H$_1$\,H$_2$ 为等腰三角形（C\,H$_1$ = C\,H$_2$），
由顶点到底边中点的连线
与底边垂直。

\textbf{15.} $\sin \widehat{\text{M\,C\,H}_1} = \dfrac{\text{M\,H}_1}{\text{C\,H}_1}
= \dfrac{a\sqrt{2}/2}{a\sqrt{3}/2} = \sqrt{2/3} \approx 0.816$。

\textbf{16.} 半角 $\approx \ang{54.7}$；键角 H--C--H $\approx
2 \times 54.7 = \ang{109.5}$。

\textbf{17.} 这正是表中四面体的 \ang{109.5}。

\textbf{18.} 钻孔时，各孔之间必须相隔 \ang{109.5}。
\end{solution}
